Unfortunately, using ML techniques has its own set of challenges. ML models typically need substantially more parameters than simple linear models to capture complex functional forms. As a result, each parameter is estimated from less information, which increases the variability of the parameter estimates relative to linear models. This creates a trade-off: even if the ML approach is unbiased on average, in practice researchers are more likely to obtain estimates that deviate substantially from the true cross-lagged effect. This trade-off is known as the bias–variance trade-off (see \citeauthor{shmueli2010explain}, \citeyear{shmueli2010explain}), which is highly relevant in this context because sample sizes in the social sciences are typically limited, with about 104 participants on average \citep{fraley2014n}. This also seems like a good time to note that decoupling covariate adjustment from model estimation leads to underestimated standard errors for the cross-lagged effects, without proper adjustment \citep{freckleton2002misuse}. This underestimation stems from the uncertainty induced by the adjustment model, which is not 'seen' by the SEM that is later fit on the residuals. This issue will be formally dealt with in the thesis. Despite these concerns, previous work using ML approaches has yielded flexible yet interpretable results in a social science context (e.g. \citeauthor{sheetal2023using}, \citeyear{sheetal2023using}), highlighting their potential to reduce misspecification bias.